\documentclass[12pt]{article}
\usepackage[a4paper,margin=1in]{geometry}
\usepackage{iftex}
\ifXeTeX
\usepackage{fontspec}
\defaultfontfeatures{Ligatures=TeX}
\setmainfont{Times New Roman}
\else
\usepackage[T1]{fontenc}
\usepackage[utf8]{inputenc}
\usepackage{newtxtext}
\usepackage{newtxmath}
\fi
\usepackage{enumitem}
\usepackage{hyperref}
\usepackage{parskip}
\usepackage{microtype}
\usepackage[normalem]{ulem}
\usepackage{xcolor}

\newcommand{\practiceid}[1]{\texttt{#1}}
\newlength{\readinggutter}
\setlength{\readinggutter}{2.8em}
\newlength{\readingfirstindent}
\setlength{\readingfirstindent}{1.5em}
\newlength{\questionblockgap}
\setlength{\questionblockgap}{0.45\baselineskip}
\newlength{\exampleblockgap}
\setlength{\exampleblockgap}{0.65\baselineskip}
\newlength{\passageparagraphgap}
\setlength{\passageparagraphgap}{0.45\baselineskip}
\newenvironment{exampleblock}{\par\vspace{0.35\baselineskip}\begingroup\raggedright\setlength{\parskip}{0pt}\setlength{\parindent}{0pt}}{\par\endgroup\vspace{\exampleblockgap}}
\newcommand{\readinglineindent}[2]{%
\par\noindent%
\hangindent=\dimexpr\readinggutter+0.9em\relax%
\hangafter=1%
\makebox[\readinggutter][r]{\scriptsize\color{gray}#1}\hspace{0.9em}\hspace*{\readingfirstindent}#2%
\par}
\newcommand{\readingline}[2]{%
\par\noindent%
\hangindent=\dimexpr\readinggutter+0.9em\relax%
\hangafter=1%
\makebox[\readinggutter][r]{\scriptsize\color{gray}#1}\hspace{0.9em}#2%
\par}

\setlength{\parindent}{0pt}
\setlength{\parskip}{0.5em}
\setlist[enumerate]{topsep=0.25em,itemsep=0.18em,parsep=0pt,partopsep=0pt}
\setlength{\emergencystretch}{2em}

\begin{document}
\section*{TOEFL Practice 01 - Tryout}
Source of truth: DOCX only.
\section*{Listening Comprehension}
\subsection*{Part A}
DIRECTIONS: In Part A you will hear short conversations between two speakers. At the end of each conversation, a third speaker will ask a question about what the first two speakers said. Each conversation and each question will be spoken only one time. Therefore, you must listen carefully to understand what each speaker says. After you hear a conversation and the question, read the four choices and select the one that is the best answer to the question the speaker asked. Then, on your answer. sheet, find the number of the question and blacken the space that corresponds to the letter for the answer you have chosen. Blacken the space completely so that the letter inside the space does not show.\par
\begin{exampleblock}
Listen to the following example.\par
On the recording, you hear:\par
(Man) Does the car need to be filled?\par
(Woman) Mary stopped at the gas station on her way home.\par
(Narrator) What does the woman mean?\par
In your test book, you will read:\par
(A) Mary bought some food.\par
(B) Mary had car trouble.\par
(C) Mary went shopping.\par
(D) Mary bought some gas.\par
From the conversation you learn that Mary stopped at the gas station on her way home. The best answer to the question "Does the car need to be filled?" is (D), "Mary bought some gas." Therefore, the correct answer is (D).\par
Now let us begin Part A with question number 1.\par
\end{exampleblock}
\begin{exampleblock}
\textbf{Transkrip rekonstruksi AI (draft; bukan resmi)}\par
Man: Mike missed class again and wants to borrow someone's notes.\par
Woman: Borrowing notes may help, but he really needs to get the material that was taught in class.\par
Narrator: What does the woman imply about Mike?\par
\end{exampleblock}
\noindent\textbf{1.}
\begin{enumerate}[label=(\Alph*),leftmargin=*,itemsep=0.2em]
\item Mike shouldn't skip classes and borrow notes.
\item Mike always borrows notes from his classmates.
\item Mike is least likely to skip classes.
\item Mike needs to get the material taught in class.
\end{enumerate}
\par
\vspace{0.25\baselineskip}
\begin{exampleblock}
\textbf{Transkrip rekonstruksi AI (draft; bukan resmi)}\par
Woman: Do you think I can still stop by the office today?\par
Man: I doubt it. They closed at five, and everyone has gone home for the day.\par
Narrator: What does the man mean?\par
\end{exampleblock}
\noindent\textbf{2.}
\begin{enumerate}[label=(\Alph*),leftmargin=*,itemsep=0.2em]
\item The office is closed for the day.
\item The officer will return in the morning.
\item There is an opening at 9 tomorrow morning.
\item The office is open until 9 o'clock.
\end{enumerate}
\par
\vspace{0.25\baselineskip}
\begin{exampleblock}
\textbf{Transkrip rekonstruksi AI (draft; bukan resmi)}\par
Man: Is something wrong with your lunch?\par
Woman: Yes, I ordered the chicken salad, but this is a green salad with no chicken at all.\par
Narrator: What does the woman mean?\par
\end{exampleblock}
\noindent\textbf{3.}
\begin{enumerate}[label=(\Alph*),leftmargin=*,itemsep=0.2em]
\item She is sorry that she ordered that salad.
\item She didn't get the salad she ordered.
\item She is sorry, but this salad is not hers.
\item This salad does not taste very good.
\end{enumerate}
\par
\vspace{0.25\baselineskip}
\begin{exampleblock}
\textbf{Transkrip rekonstruksi AI (draft; bukan resmi)}\par
Woman: Did Mr. Calvert say when he is planning to return?\par
Man: He mentioned sometime next week, but he did not give an exact day.\par
Narrator: What did the woman ask?\par
\end{exampleblock}
\noindent\textbf{4.}
\begin{enumerate}[label=(\Alph*),leftmargin=*,itemsep=0.2em]
\item Did Mr. Calvert say something about his back?
\item Did Mr. Calvert say when he is planning to return?
\item Mr. Calvert forgot it's time for him to come back.
\item Mr. Calvert said he isn't going to come back.
\end{enumerate}
\par
\vspace{0.25\baselineskip}
\begin{exampleblock}
\textbf{Transkrip rekonstruksi AI (draft; bukan resmi)}\par
Man: The book I need is checked out. What should I do?\par
Woman: Leave your name at the desk. The library will notify you when the book arrives.\par
Narrator: What does the woman say?\par
\end{exampleblock}
\noindent\textbf{5.}
\begin{enumerate}[label=(\Alph*),leftmargin=*,itemsep=0.2em]
\item The library sells books and postcards.
\item Postcards are available to library users.
\item You need to have a card to have your books sent.
\item The library will let you know when the book arrives.
\end{enumerate}
\par
\vspace{0.25\baselineskip}
\begin{exampleblock}
\textbf{Transkrip rekonstruksi AI (draft; bukan resmi)}\par
Woman: Did you cook dinner last night?\par
Man: No, I was too tired, so I just went out to eat.\par
Narrator: What does the man mean?\par
\end{exampleblock}
\noindent\textbf{6.}
\begin{enumerate}[label=(\Alph*),leftmargin=*,itemsep=0.2em]
\item He does not know how to cook.
\item He had dinner at a restaurant.
\item He doesn't like to stand while cooking.
\item He went shopping and then had dinner.
\end{enumerate}
\par
\vspace{0.25\baselineskip}
\begin{exampleblock}
\textbf{Transkrip rekonstruksi AI (draft; bukan resmi)}\par
Man: I want to bring my bike to campus. Do I need anything?\par
Woman: Yes, bicycle permits are available at the student services office next door.\par
Narrator: What does the woman mean?\par
\end{exampleblock}
\noindent\textbf{7.}
\begin{enumerate}[label=(\Alph*),leftmargin=*,itemsep=0.2em]
\item Students are permitted to use bikes.
\item Students with bikes are counted regularly.
\item Permits for bikes are available.
\item Bikes can be purchased next door.
\end{enumerate}
\par
\vspace{0.25\baselineskip}
\begin{exampleblock}
\textbf{Transkrip rekonstruksi AI (draft; bukan resmi)}\par
Woman: When is the paper due?\par
Man: I have no idea. They still have not told us when to turn it in.\par
Narrator: What does the man mean?\par
\end{exampleblock}
\noindent\textbf{8.}
\begin{enumerate}[label=(\Alph*),leftmargin=*,itemsep=0.2em]
\item They haven't told him when to turn in the paper.
\item They told him when the assignment is due, but he forgot.
\item He wishes they would tell him where to pay the dues.
\item He wanted to tell them when to turn in the assignment.
\end{enumerate}
\par
\vspace{0.25\baselineskip}
\begin{exampleblock}
\textbf{Transkrip rekonstruksi AI (draft; bukan resmi)}\par
Man: What happened in the driveway?\par
Woman: The driver backed up too far and hit the garbage can by the garage.\par
Narrator: What does the woman say?\par
\end{exampleblock}
\noindent\textbf{9.}
\begin{enumerate}[label=(\Alph*),leftmargin=*,itemsep=0.2em]
\item The garbage can be stored in the garage.
\item The driver is parking in the garage.
\item The driver crashed into the garage.
\item The garbage can was near the garage.
\end{enumerate}
\par
\vspace{0.25\baselineskip}
\begin{exampleblock}
\textbf{Transkrip rekonstruksi AI (draft; bukan resmi)}\par
Woman: Joan was sitting in the back of the room during the speech.\par
Man: Then I am not surprised she missed so much of it. The speaker was not loud enough.\par
Narrator: What does the man imply?\par
\end{exampleblock}
\noindent\textbf{10.}
\begin{enumerate}[label=(\Alph*),leftmargin=*,itemsep=0.2em]
\item Joan didn't hear the first half of the speech.
\item Joan heard nothing from half of the speakers.
\item Half of the people in the room did not hear the speaker.
\item The speaker did not speak loudly enough.
\end{enumerate}
\par
\vspace{0.25\baselineskip}
\begin{exampleblock}
\textbf{Transkrip rekonstruksi AI (draft; bukan resmi)}\par
Man: Can I make calls outside campus from my room phone now?\par
Woman: Not yet. You cannot make off-campus calls until the phone service is connected.\par
Narrator: What does the woman mean?\par
\end{exampleblock}
\noindent\textbf{11.}
\begin{enumerate}[label=(\Alph*),leftmargin=*,itemsep=0.2em]
\item On campus, your phone can be hooked up only once.
\item Campus calls can't be made from your phone.
\item You cannot call off campus until your phone is connected.
\item You can't call collect unless you live on campus.
\end{enumerate}
\par
\vspace{0.25\baselineskip}
\begin{exampleblock}
\textbf{Transkrip rekonstruksi AI (draft; bukan resmi)}\par
Woman: I stopped by the drugstore to buy toothpaste.\par
Man: Were you able to get some?\par
Woman: No, they were completely out.\par
Narrator: What does the woman mean?\par
\end{exampleblock}
\noindent\textbf{12.}
\begin{enumerate}[label=(\Alph*),leftmargin=*,itemsep=0.2em]
\item The drug store sells toothpaste.
\item Drugstores used to sell toothpaste.
\item The drugstore does not have toothpaste.
\item Toothpaste and drugs are sold everywhere.
\end{enumerate}
\par
\vspace{0.25\baselineskip}
\begin{exampleblock}
\textbf{Transkrip rekonstruksi AI (draft; bukan resmi)}\par
Man: That truck just passed another car on the right.\par
Woman: That driver could get a ticket. Passing on the right is against the law here.\par
Narrator: What does the woman mean?\par
\end{exampleblock}
\noindent\textbf{13.}
\begin{enumerate}[label=(\Alph*),leftmargin=*,itemsep=0.2em]
\item The truck on your right has no ticket.
\item If you buy a ticket, you can pass the driving test.
\item Passing on the right is against the law.
\item Trucks get tickets for passing, but cars don't.
\end{enumerate}
\par
\vspace{0.25\baselineskip}
\begin{exampleblock}
\textbf{Transkrip rekonstruksi AI (draft; bukan resmi)}\par
Woman: Why did he choose the local college?\par
Man: He did not want to go far from home.\par
Narrator: What does the man mean?\par
\end{exampleblock}
\noindent\textbf{14.}
\begin{enumerate}[label=(\Alph*),leftmargin=*,itemsep=0.2em]
\item He didn't travel to college.
\item He didn't attend college.
\item He didn't go to school.
\item He didn't want to go far.
\end{enumerate}
\par
\vspace{0.25\baselineskip}
\begin{exampleblock}
\textbf{Transkrip rekonstruksi AI (draft; bukan resmi)}\par
Man: Why are so many forms filled out incorrectly?\par
Woman: Most students simply ignore the directions on the paperwork.\par
Narrator: What does the woman mean?\par
\end{exampleblock}
\noindent\textbf{15.}
\begin{enumerate}[label=(\Alph*),leftmargin=*,itemsep=0.2em]
\item Students don't read their instructor's forms.
\item If students can't read, they can't fill out forms.
\item Most students ignore the directions on paperwork.
\item Most forms for students don't have instructions.
\end{enumerate}
\par
\vspace{0.25\baselineskip}
\begin{exampleblock}
\textbf{Transkrip rekonstruksi AI (draft; bukan resmi)}\par
Woman: How much farther is the movie theater?\par
Man: It is just around the corner. Turn left at the next street.\par
Narrator: What does the man mean?\par
\end{exampleblock}
\noindent\textbf{16.}
\begin{enumerate}[label=(\Alph*),leftmargin=*,itemsep=0.2em]
\item The movie theater is around the corner.
\item The movie theater is to your left.
\item The movie theater is a mile away.
\item The next turn is by the movie theater.
\end{enumerate}
\par
\vspace{0.25\baselineskip}
\begin{exampleblock}
\textbf{Transkrip rekonstruksi AI (draft; bukan resmi)}\par
Man: I heard Julie is taking five courses this term.\par
Woman: Yes, she is carrying a full course load.\par
Narrator: What does the woman mean?\par
\end{exampleblock}
\noindent\textbf{17.}
\begin{enumerate}[label=(\Alph*),leftmargin=*,itemsep=0.2em]
\item Julie's courses are full and closed.
\item Julie is carrying a bag full of books.
\item Julie can't carry all these courses.
\item Julie is enrolled as a full-time student.
\end{enumerate}
\par
\vspace{0.25\baselineskip}
\begin{exampleblock}
\textbf{Transkrip rekonstruksi AI (draft; bukan resmi)}\par
Woman: How was the field trip?\par
Man: It was interesting, but not many students went.\par
Narrator: What does the man mean?\par
\end{exampleblock}
\noindent\textbf{18.}
\begin{enumerate}[label=(\Alph*),leftmargin=*,itemsep=0.2em]
\item Everyone who went on the trip fell.
\item The trip was not well attended.
\item The trip was hard to make in one day.
\item Many students tripped and fell.
\end{enumerate}
\par
\vspace{0.25\baselineskip}
\begin{exampleblock}
\textbf{Transkrip rekonstruksi AI (draft; bukan resmi)}\par
Woman: Did either swim team win the meet?\par
Man: No, the men's and women's teams both ended in a tie.\par
Narrator: What does the man mean?\par
\end{exampleblock}
\noindent\textbf{19.}
\begin{enumerate}[label=(\Alph*),leftmargin=*,itemsep=0.2em]
\item Neither team won.
\item Men swim better than women.
\item Did the men's team win?
\item Is Tim the head of the team?
\end{enumerate}
\par
\vspace{0.25\baselineskip}
\begin{exampleblock}
\textbf{Transkrip rekonstruksi AI (draft; bukan resmi)}\par
Man: Has she had time to read the report yet?\par
Woman: I do not think so. She has been in meetings all morning.\par
Narrator: What did the man want to know?\par
\end{exampleblock}
\noindent\textbf{20.}
\begin{enumerate}[label=(\Alph*),leftmargin=*,itemsep=0.2em]
\item She had a good chance to see the report.
\item Chances are she was looking at a reporter.
\item Has this report been changed yet?
\item Has she had time to read the report?
\end{enumerate}
\par
\vspace{0.25\baselineskip}
\begin{exampleblock}
\textbf{Transkrip rekonstruksi AI (draft; bukan resmi)}\par
Woman: It was sunny this morning, but now it looks like rain.\par
Man: That is typical here. The weather is impossible to predict.\par
Narrator: What does the man mean?\par
\end{exampleblock}
\noindent\textbf{21.}
\begin{enumerate}[label=(\Alph*),leftmargin=*,itemsep=0.2em]
\item The weather is unpredictable.
\item Rain is predicted for today.
\item She heard the weather report.
\item This year, it's been raining often.
\end{enumerate}
\par
\vspace{0.25\baselineskip}
\begin{exampleblock}
\textbf{Transkrip rekonstruksi AI (draft; bukan resmi)}\par
Man: How long will it take to fix my brakes?\par
Woman: The mechanic can check them this afternoon, but the repair may take until tomorrow.\par
Narrator: Where does this conversation probably take place?\par
\end{exampleblock}
\noindent\textbf{22.}
\begin{enumerate}[label=(\Alph*),leftmargin=*,itemsep=0.2em]
\item At a police station
\item At a doctor's office
\item In a car repair shop
\item In an insurance agency
\end{enumerate}
\par
\vspace{0.25\baselineskip}
\begin{exampleblock}
\textbf{Transkrip rekonstruksi AI (draft; bukan resmi)}\par
Woman: Can your company mow the front yard and trim the hedges this week?\par
Man: Yes, our lawn crew can come by on Friday.\par
Narrator: What kind of business is being discussed?\par
\end{exampleblock}
\noindent\textbf{23.}
\begin{enumerate}[label=(\Alph*),leftmargin=*,itemsep=0.2em]
\item A house-cleaning business
\item A plumbing company
\item A lawn service
\item An electrical supply store
\end{enumerate}
\par
\vspace{0.25\baselineskip}
\begin{exampleblock}
\textbf{Transkrip rekonstruksi AI (draft; bukan resmi)}\par
Man: What are you going to do after your last exam?\par
Woman: I think I will take a break before starting anything else.\par
Narrator: What does the woman plan to do?\par
\end{exampleblock}
\noindent\textbf{24.}
\begin{enumerate}[label=(\Alph*),leftmargin=*,itemsep=0.2em]
\item Finish her classes
\item Look for a job
\item Write a paper
\item Take a break
\end{enumerate}
\par
\vspace{0.25\baselineskip}
\begin{exampleblock}
\textbf{Transkrip rekonstruksi AI (draft; bukan resmi)}\par
Woman: I am hungry, but I do not want a big dinner.\par
Man: There is a small cafe across the street. Let's go there.\par
Narrator: Where are they probably going?\par
\end{exampleblock}
\noindent\textbf{25.}
\begin{enumerate}[label=(\Alph*),leftmargin=*,itemsep=0.2em]
\item Over to the next street
\item For dinner at a shop
\item To a cafe
\item For a walk
\end{enumerate}
\par
\vspace{0.25\baselineskip}
\begin{exampleblock}
\textbf{Transkrip rekonstruksi AI (draft; bukan resmi)}\par
Woman: The water heater in my new apartment had to be replaced.\par
Man: A new heater must have been expensive.\par
Narrator: What does the man imply?\par
\end{exampleblock}
\noindent\textbf{26.}
\begin{enumerate}[label=(\Alph*),leftmargin=*,itemsep=0.2em]
\item The move in hot weather was costly.
\item The woman must have moved too soon.
\item A new heater must have been expensive.
\item The water heater never worked well.
\end{enumerate}
\par
\vspace{0.25\baselineskip}
\begin{exampleblock}
\textbf{Transkrip rekonstruksi AI (draft; bukan resmi)}\par
Man: The department is replacing all the old computers this month.\par
Woman: It is about time. They have needed new ones for a while.\par
Narrator: What does the woman mean?\par
\end{exampleblock}
\noindent\textbf{27.}
\begin{enumerate}[label=(\Alph*),leftmargin=*,itemsep=0.2em]
\item They don't need new computers at this time.
\item They don't have much time to work with new computers.
\item The computers should have been replaced much earlier.
\item The timing for replacing computers is about right.
\end{enumerate}
\par
\vspace{0.25\baselineskip}
\begin{exampleblock}
\textbf{Transkrip rekonstruksi AI (draft; bukan resmi)}\par
Man: Could you come by tomorrow to help me with the application?\par
Woman: Tomorrow will not work, but I can come another day this week.\par
Narrator: What does the woman mean?\par
\end{exampleblock}
\noindent\textbf{28.}
\begin{enumerate}[label=(\Alph*),leftmargin=*,itemsep=0.2em]
\item She can come on another day, but not tomorrow.
\item She can see the man in three days.
\item She is thirsty and needs some water.
\item She doesn't know how to help the man.
\end{enumerate}
\par
\vspace{0.25\baselineskip}
\begin{exampleblock}
\textbf{Transkrip rekonstruksi AI (draft; bukan resmi)}\par
Man: Do you have change for a dollar? The vending machine will not take my bill.\par
Woman: I think I have some coins in my bag.\par
Narrator: What does the man need?\par
\end{exampleblock}
\noindent\textbf{29.}
\begin{enumerate}[label=(\Alph*),leftmargin=*,itemsep=0.2em]
\item The man needs change for a vending machine.
\item The man will break expensive equipment.
\item The man shouldn't spend his money at the bar.
\item The man shouldn't give her candy.
\end{enumerate}
\par
\vspace{0.25\baselineskip}
\begin{exampleblock}
\textbf{Transkrip rekonstruksi AI (draft; bukan resmi)}\par
Woman: Could you push the cart while I get the apples?\par
Man: Sure. I will wait here by the dairy section.\par
Narrator: Where are the speakers?\par
\end{exampleblock}
\noindent\textbf{30.}
\begin{enumerate}[label=(\Alph*),leftmargin=*,itemsep=0.2em]
\item On the street
\item In the car
\item In a warehouse
\item At a grocery store.
\end{enumerate}
\par
\vspace{0.25\baselineskip}
\subsection*{Part B}
DIRECTIONS: In this part of the test, you will hear longer conversations. After each conversation, you will hear several questions. The conversations and questions will not be repeated. After you hear a question, read the four possible answers in. your test book and select. The best answer. Then, on your answer sheet, find the number of the question and fill in the space that corresponds to the letter of the answer you have chosen. Remember, you are not allowed to take notes or write in your test book.\par
\begin{exampleblock}
Listen to the following example:\par
You will hear:\par
You will read:\par
(A) He has changed jobs.\par
(B) He has two children.\par
(C) He has two jobs.\par
(D) He is looking for a job.\par
From the conversation you learn that Tom has taken an additional job. The best answer to the question "Why is Tom tired?" is (C), "He has two jobs." Therefore the correct answer is (C).\par
\end{exampleblock}
\begin{exampleblock}
\textbf{Transkrip rekonstruksi AI (draft; bukan resmi)}\par
Woman: Excuse me, do you know anything about glue? The handle on my purse ripped this morning, and I am not sure what to use.\par
Man: It depends on the material. Ordinary white glue works well for paper and light wood, but not for leather or plastic.\par
Woman: I was thinking about rubber cement. Would that hold it?\par
Man: Rubber cement is flexible, but it is not strong enough for something that gets pulled a lot, like a purse handle. For that, epoxy would probably be better.\par
Woman: Epoxy? I thought that was only for heavy repairs.\par
Man: Not always. I used epoxy last month when my daughter's shoe started coming apart, and it held surprisingly well.\par
Woman: So the shoe was just an example of how epoxy can be used?\par
Man: Exactly. If you want, I can show you which kind to buy and help you fix the purse.\par
Woman: That would be great. I really appreciate it.\par
\end{exampleblock}
\noindent\textbf{31.}
\begin{enumerate}[label=(\Alph*),leftmargin=*,itemsep=0.2em]
\item Children's shoes
\item Business trips
\item Different types of glue
\item Various types of goods
\end{enumerate}
\par
\vspace{0.25\baselineskip}
\noindent\textbf{32.}
\begin{enumerate}[label=(\Alph*),leftmargin=*,itemsep=0.2em]
\item She had to leave early in the morning.
\item She needed to call her friend.
\item Her purse had ripped.
\item Her office didn't supply the hardware.
\end{enumerate}
\par
\vspace{0.25\baselineskip}
\noindent\textbf{33.}
\begin{enumerate}[label=(\Alph*),leftmargin=*,itemsep=0.2em]
\item Come to her house
\item Glue her purse
\item Suggest an adhesive
\item Go to a store to buy glue
\end{enumerate}
\par
\vspace{0.25\baselineskip}
\noindent\textbf{34.}
\begin{enumerate}[label=(\Alph*),leftmargin=*,itemsep=0.2em]
\item To show that he is a good father
\item To display his knowledge
\item As an indication of problems with cement
\item As an example of using epoxy
\end{enumerate}
\par
\vspace{0.25\baselineskip}
\noindent\textbf{35.}
\begin{enumerate}[label=(\Alph*),leftmargin=*,itemsep=0.2em]
\item He doesn't like to be bothered.
\item He wants to help the woman.
\item He wants the woman to like him.
\item He isn't interested in women's purses.
\end{enumerate}
\par
\vspace{0.25\baselineskip}
\begin{exampleblock}
\textbf{Transkrip rekonstruksi AI (draft; bukan resmi)}\par
Man: Is this your first time in Portland?\par
Woman: Yes. I came down from Seattle for the conference. I heard a radio report about the event before I left.\par
Man: That is interesting. I live here, but I do not always follow national newspapers. I usually pay attention to regional news.\par
Woman: That makes sense. Local news probably covers this convention better than the national papers do.\par
Man: Exactly. The Portland stations have been talking about it all week.\par
\end{exampleblock}
\noindent\textbf{36.}
\begin{enumerate}[label=(\Alph*),leftmargin=*,itemsep=0.2em]
\item New York
\item Seattle
\item Vancouver
\item Portland
\end{enumerate}
\par
\vspace{0.25\baselineskip}
\noindent\textbf{37.}
\begin{enumerate}[label=(\Alph*),leftmargin=*,itemsep=0.2em]
\item She gets the New York Times.
\item She reads about it in the Seattle Post.
\item She listens to the news on the radio.
\item She watches TV and goes to the movies.
\end{enumerate}
\par
\vspace{0.25\baselineskip}
\noindent\textbf{38.}
\begin{enumerate}[label=(\Alph*),leftmargin=*,itemsep=0.2em]
\item He reads the newspaper and listens to the radio.
\item He travels to Portland to attend conventions.
\item Usually, he pays attention to regional news.
\item He talks to tourists who come to visit the city.
\end{enumerate}
\par
\vspace{0.25\baselineskip}
\subsection*{Part C}
DIRECTIONS: In Part C you will hear short lectures and extended conversations. At the end of each, you will be asked several questions. Each lecture or conversation and each question will be spoken only one time. For this reason, you must listen carefully to understand what each speaker says. After you hear a question, read the four choices and select the one that best answers the question the speaker asked. Then, on your answer sheet, find the number of the question and blacken the space that contains the letter for the answer you have chosen. Answer all questions according to what is stated or implied in the lecture or conversation.\par
\begin{exampleblock}
Listen to this sample talk.\par
You will hear:\par
Now listen, to the following example.\par
You will hear:\par
You will read:\par
(A) By cars and carriages\par
(B) By bicycles, trains, and carriages\par
(C) On foot and by boat\par
(D) On board ships and trains\par
The best answer to the question "According to the speaker, how did people travel before the invention of the automobile?" is (B), "By bicycles, trains, and carriages." Therefore, the correct answer is (B). You are not allowed to make notes during the test.\par
\end{exampleblock}
\begin{exampleblock}
\textbf{Transkrip rekonstruksi AI (draft; bukan resmi)}\par
Air transportation is used not only for passengers but also for cargo. In large airports, cargo handling is a separate and highly organized operation. Air cargo usually includes two major categories: mail and freight. Mail is often processed through postal facilities located in or near airports, while freight is routed to separate cargo terminals.\par
Once freight reaches the airport, it is weighed, labeled, sorted, and sent to the correct terminal area. Some freight travels on passenger aircraft, but other freight is shipped according to special cargo schedules. This allows airlines to move goods quickly even when passenger flights are not available or when the cargo requires special handling.\par
Cargo workers also coordinate with customs officers, postal workers, and airline staff to make sure shipments are loaded on the correct aircraft. Because timing is important, much of this work takes place in airport facilities designed specifically for cargo movement.\par
\end{exampleblock}
\noindent\textbf{39.}
\begin{enumerate}[label=(\Alph*),leftmargin=*,itemsep=0.2em]
\item How mail is processed in airports
\item How quickly goods can be delivered
\item How big and small airports work
\item How cargo is handled for shipment
\end{enumerate}
\par
\vspace{0.25\baselineskip}
\noindent\textbf{40.}
\begin{enumerate}[label=(\Alph*),leftmargin=*,itemsep=0.2em]
\item Mail and baggage
\item Mail and freight
\item Passengers and luggage
\item Baggage and freight
\end{enumerate}
\par
\vspace{0.25\baselineskip}
\noindent\textbf{41.}
\begin{enumerate}[label=(\Alph*),leftmargin=*,itemsep=0.2em]
\item It is routed to separate cargo terminals.
\item It is separated from mail by machine.
\item It is loaded when passengers board the planes.
\item It is shipped to smaller airports.
\end{enumerate}
\par
\vspace{0.25\baselineskip}
\noindent\textbf{42.}
\begin{enumerate}[label=(\Alph*),leftmargin=*,itemsep=0.2em]
\item Together with passengers
\item Together with baggage
\item According to special schedules
\item Exclusively to its destination
\end{enumerate}
\par
\vspace{0.25\baselineskip}
\noindent\textbf{43.}
\begin{enumerate}[label=(\Alph*),leftmargin=*,itemsep=0.2em]
\item In postal offices
\item In airports
\item On aircraft
\item In commercial outlets
\end{enumerate}
\par
\vspace{0.25\baselineskip}
\begin{exampleblock}
\textbf{Transkrip rekonstruksi AI (draft; bukan resmi)}\par
Today we will look at bananas as an agricultural crop. Bananas are not manufactured goods, and they are not related to passenger travel or automobiles. They belong to the category of crops, and they are one of the most widely consumed fruits in the world.\par
There are thousands of varieties of bananas, though only a few are commonly sold in stores. Banana plants produce flowers that may be green, white, or cream-colored. The fruit itself can be described in several ways depending on the variety. Some bananas are soft when ripe, while others remain firm and are often cooked before being eaten.\par
The most familiar type is the sweet banana, which is usually eaten fresh. Other kinds may be cooked or used in different food products, but dessert bananas are generally eaten as fruit. Banana plants require warmth, moisture, and strong root systems. Their extensive roots help them take in water and nutrients from the soil.\par
Commercial banana production is especially important in tropical regions, including parts of South America. From there, bananas are exported to markets around the world.\par
\end{exampleblock}
\noindent\textbf{44.}
\begin{enumerate}[label=(\Alph*),leftmargin=*,itemsep=0.2em]
\item Manufactured goods
\item Passengers and baggage
\item Livestock and crops
\item Cars and automobile mechanics
\end{enumerate}
\par
\vspace{0.25\baselineskip}
\noindent\textbf{45.}
\begin{enumerate}[label=(\Alph*),leftmargin=*,itemsep=0.2em]
\item Thousands
\item A few
\item Thirty
\item Sixty
\end{enumerate}
\par
\vspace{0.25\baselineskip}
\noindent\textbf{46.}
\begin{enumerate}[label=(\Alph*),leftmargin=*,itemsep=0.2em]
\item Green, yellow, and red
\item Green, white, or cream.
\item Yellow, white, or red
\item White, cream, and yellow.
\end{enumerate}
\par
\vspace{0.25\baselineskip}
\noindent\textbf{47.}
\begin{enumerate}[label=(\Alph*),leftmargin=*,itemsep=0.2em]
\item As soft or firm
\item As sweet or sour
\item As fresh or sour
\item As large or small
\end{enumerate}
\par
\vspace{0.25\baselineskip}
\noindent\textbf{48.}
\begin{enumerate}[label=(\Alph*),leftmargin=*,itemsep=0.2em]
\item They are eaten.
\item They are pressed for juice.
\item They are made into sauce.
\item They are processed as jam.
\end{enumerate}
\par
\vspace{0.25\baselineskip}
\noindent\textbf{49.}
\begin{enumerate}[label=(\Alph*),leftmargin=*,itemsep=0.2em]
\item Cold weather
\item Prolonged watering
\item Systematic consumption
\item Extensive roots
\end{enumerate}
\par
\vspace{0.25\baselineskip}
\noindent\textbf{50.}
\begin{enumerate}[label=(\Alph*),leftmargin=*,itemsep=0.2em]
\item Africa
\item South America
\item Australia
\item Eurasia
\end{enumerate}
\par
\vspace{0.25\baselineskip}
\section*{Structure and Written Expression}
\subsection*{Sentence Completion}
This section is designed to test your ability to recognize language structures that are appropriate in standard written English. The questions in this section belong to two types, each of which has special directions.\par
DIRECTIONS: Questions 1-15 are partial sentences. Below each sentence you will see four words or phrases, marked (A), (B), (C), and (D). Select the one word or phrase that best completes the sentence. Then, on your answer sheet, find the number of the question and fill in the space that contains the letter for the answer you have chosen. Fill in the space completely.\par
\begin{exampleblock}
\textbf{EXAMPLE I}\par
\noindent Drying flowers is the best way......them.
\begin{enumerate}[label=(\Alph*),leftmargin=*,itemsep=0.2em]
\item to preserve
\item by preserving
\item preserve
\item preserved
\end{enumerate}
\vspace{0.25\baselineskip}
The sentence should state, "Drying flowers is the best way to preserve them." Therefore, the correct answer is (A).\par
\end{exampleblock}
\begin{exampleblock}
\textbf{EXAMPLE II}\par
\noindent Many American universities......as small, private colleges.
\begin{enumerate}[label=(\Alph*),leftmargin=*,itemsep=0.2em]
\item begun
\item beginning
\item began
\item for the beginning
\end{enumerate}
\vspace{0.25\baselineskip}
The sentence should state, "Many American universities began as small, private colleges." Therefore, the correct answer is (C). After you read the directions, begin work on the questions.\par
\end{exampleblock}
\noindent\textbf{1.} Fort Niagara was built by the French in 1726 on land.......the Seneca Indians.
\begin{enumerate}[label=(\Alph*),leftmargin=*,itemsep=0.2em]
\item they buy from
\item bought from
\item buying from
\item was bought from
\end{enumerate}
\par
\vspace{0.25\baselineskip}
\noindent\textbf{2.} Soil texture depends on the proportions of clay and sand particles,......soil porosity.
\begin{enumerate}[label=(\Alph*),leftmargin=*,itemsep=0.2em]
\item both alter
\item which alter
\item where altered
\item although altered
\end{enumerate}
\par
\vspace{0.25\baselineskip}
\noindent\textbf{3.} The writers of the realist movement embraced the notion that art should depict life......
\begin{enumerate}[label=(\Alph*),leftmargin=*,itemsep=0.2em]
\item accurately and objectively
\item accuracy and objectivity
\item accurate and objective
\item accurate objectivity
\end{enumerate}
\par
\vspace{0.25\baselineskip}
\noindent\textbf{4.} A ratio is a comparison of......whole or a part to another part.
\begin{enumerate}[label=(\Alph*),leftmargin=*,itemsep=0.2em]
\item part to the
\item a part to
\item a part to the
\item the part to the
\end{enumerate}
\par
\vspace{0.25\baselineskip}
\noindent\textbf{5.} The bones of the elderly are more prone to fractures and splintering......of young people.
\begin{enumerate}[label=(\Alph*),leftmargin=*,itemsep=0.2em]
\item than that
\item than those
\item those than
\item that than
\end{enumerate}
\par
\vspace{0.25\baselineskip}
\noindent\textbf{6.} English and Scottish settlers......Belfast as a trading post in 1613.
\begin{enumerate}[label=(\Alph*),leftmargin=*,itemsep=0.2em]
\item they established
\item established themselves
\item established
\item establishing
\end{enumerate}
\par
\vspace{0.25\baselineskip}
\noindent\textbf{7.} The formulation of economic policies necessitates meticulous consideration......large segments of the population.
\begin{enumerate}[label=(\Alph*),leftmargin=*,itemsep=0.2em]
\item because they affect
\item they are affected because
\item affect them because
\item because affecting them
\end{enumerate}
\par
\vspace{0.25\baselineskip}
\noindent\textbf{8.} Only......feathered creatures inhabit the Arctic region year round.
\begin{enumerate}[label=(\Alph*),leftmargin=*,itemsep=0.2em]
\item fewer
\item fewer than
\item as few as
\item a few
\end{enumerate}
\par
\vspace{0.25\baselineskip}
\noindent\textbf{9.} Before Richard Bennett accepted the appointment as the prime minister of Canada in 1930, he......as a lawyer.
\begin{enumerate}[label=(\Alph*),leftmargin=*,itemsep=0.2em]
\item had achieved a successful
\item had been achieved successfully
\item has achieved success
\item had achieved success
\end{enumerate}
\par
\vspace{0.25\baselineskip}
\noindent\textbf{10.} Gardeners transplant bushes and flowers by moving them from one place to......
\begin{enumerate}[label=(\Alph*),leftmargin=*,itemsep=0.2em]
\item other
\item others
\item another
\item each other
\end{enumerate}
\par
\vspace{0.25\baselineskip}
\noindent\textbf{11.} Museums of natural history are ordinarily.......by special interest groups created for that purpose.
\begin{enumerate}[label=(\Alph*),leftmargin=*,itemsep=0.2em]
\item owned and operated
\item they own and operate
\item owning and operating
\item the owner operates
\end{enumerate}
\par
\vspace{0.25\baselineskip}
\noindent\textbf{12.} A surge in the level of stress......the recurrence of nightmares.
\begin{enumerate}[label=(\Alph*),leftmargin=*,itemsep=0.2em]
\item apparent increase
\item apparently increase
\item apparently increases
\item apparent increases
\end{enumerate}
\par
\vspace{0.25\baselineskip}
\noindent\textbf{13.} Each bowler......in each frame, unless a strike is bowled.
\begin{enumerate}[label=(\Alph*),leftmargin=*,itemsep=0.2em]
\item rolling the ball twice
\item the ball is rolled twice
\item rolls the ball twice
\item the ball rolls twice
\end{enumerate}
\par
\vspace{0.25\baselineskip}
\noindent\textbf{14.} William Hearst had five sons,......eventually became executives in the Hearst newspaper conglomerate.
\begin{enumerate}[label=(\Alph*),leftmargin=*,itemsep=0.2em]
\item all of them
\item of them all
\item all of whom
\item who of all
\end{enumerate}
\par
\vspace{0.25\baselineskip}
\noindent\textbf{15.} An axiomatic assumption in physics holds that all matter has kinetic energy......motion and mass.
\begin{enumerate}[label=(\Alph*),leftmargin=*,itemsep=0.2em]
\item because its
\item because of its
\item because it is
\item because of it
\end{enumerate}
\par
\vspace{0.25\baselineskip}
\subsection*{Error Identification}
DIRECTIONS: In questions 16-40 every sentence has four words or phrases that are underlined. The four underlined portions of each sentence are marked (A), (B), (C), and (D). Identify the one word or phrase that makes the sentence incorrect. Then, on your answer sheet, find the number of the question and fill in the space that contains the letter for the answer you have chosen. Fill in the space completely.\par
\begin{exampleblock}
\textbf{EXAMPLE I}\par
Christopher Columbus \uline{has sailed}(A) from Europe in 1492 and \uline{discovered a}(B) \uline{new land}(C) he \uline{thought to}(D) be India.\par
The sentence should state, "Christopher Columbus sailed from Europe in 1492 and discovered a new land he thought to be India." Therefore, you should choose answer (A).\par
\end{exampleblock}
\begin{exampleblock}
\textbf{EXAMPLE II}\par
\uline{As the roles}(A) of people in society change, \uline{so does}(B) the \uline{rules of}(C) conduct in certain \uline{situations}(D).\par
The sentence should state, "As the roles of people in society change, so do the rules of\par
conduct in certain situations." Therefore, you should choose answer (B).\par
\end{exampleblock}
\noindent\textbf{16.} \uline{In summer}(A), warm southern \uline{air carries}(B) \uline{moist}(C) north to the eastern and \uline{central}(D) United States.
\par
\vspace{\questionblockgap}
\noindent\textbf{17.} Billie Holiday became \uline{recognized}(A) as the most \uline{innovative}(B) jazz singer of her \uline{day}(C) and was \uline{admiration}(D) for her vocal range.
\par
\vspace{\questionblockgap}
\noindent\textbf{18.} To raise livestock \uline{successfully}(A), farmers must \uline{selecting}(B) cattle for \uline{breeding}(C) and \uline{apply}(D) a dietary regimen.
\par
\vspace{\questionblockgap}
\noindent\textbf{19.} In the 1960s, \uline{urban renewal}(A) projects \uline{cleared land}(B) for \uline{commerce}(C) and \uline{offices building.}(D)
\par
\vspace{\questionblockgap}
\noindent\textbf{20.} In 1868, Sioux leaders signed a \uline{treaty}(A) \uline{preventing}(B) whites from \uline{traveling}(C) through the Sioux \uline{territorial.}(D)
\par
\vspace{\questionblockgap}
\noindent\textbf{21.} \uline{A number}(A) multiplied \uline{by zero}(B) is zero, and a number \uline{multiplied}(C) by one is the same \uline{as number}(D).
\par
\vspace{\questionblockgap}
\noindent\textbf{22.} Muscles \uline{aids}(A) in \uline{attaching}(B) portions of the skeleton to one \uline{another}(C) and ultimately \uline{shape}(D) the torso.
\par
\vspace{\questionblockgap}
\noindent\textbf{23.} Thomas More, who \uline{fell into}(A) disfavor with the king, \uline{was a}(B) great English author, \uline{statesman}(C), and \uline{scholars.}(D)
\par
\vspace{\questionblockgap}
\noindent\textbf{24.} \uline{The first}(A) microprocessors \uline{were fabricated}(B) in 1971 \uline{for installation}(C) \uline{in handhold}(D) calculators.
\par
\vspace{\questionblockgap}
\noindent\textbf{25.} If autistic \uline{children}(A) form an \uline{attachment}(B), it \uline{predominantly was}(C) to \uline{inanimate objects}(D).
\par
\vspace{\questionblockgap}
\noindent\textbf{26.} Technology \uline{is define}(A) as \uline{the tools}(B), skills, and methods \uline{that are}(C) necessary \uline{to produce}(D) goods.
\par
\vspace{\questionblockgap}
\noindent\textbf{27.} Fruit flies do \uline{not have}(A) to \uline{leap to}(B) take off \uline{because of}(C) they become airborne solely by \uline{wing movement}(D).
\par
\vspace{\questionblockgap}
\noindent\textbf{28.} Historians postulate \uline{that}(A) Eskimos \uline{migrated}(B) from Alaska to Greenland \uline{in two}(C) \uline{greater}(D) movements.
\par
\vspace{\questionblockgap}
\noindent\textbf{29.} \uline{Electric}(A) wires \uline{carry current}(B) for \uline{lighting}(C) and outlets \uline{designing}(D) for household appliances.
\par
\vspace{\questionblockgap}
\noindent\textbf{30.} Troops \uline{housing}(A) in Fort Bliss, Texas, \uline{train}(B) to \uline{operate}(C) aircraft \uline{equipment}(D) and artillery.
\par
\vspace{\questionblockgap}
\noindent\textbf{31.} Charles Kettering \uline{patented}(A) the first \uline{success}(B) spark-\uline{based}(C) starter for \uline{automotive}(D) vehicles in 1911.
\par
\vspace{\questionblockgap}
\noindent\textbf{32.} During the 1700s, public concerts \uline{proliferated}(A) when composers \uline{wrote}(B) music for \uline{their}(C)audiences \uline{enjoying.}(D)
\par
\vspace{\questionblockgap}
\noindent\textbf{33.} The \uline{philosophers}(A) and artists of ancient Greece and Rome \uline{emphasized the}(B) study of \uline{human as}(C) \uline{fundamental}(D) to their doctrine.
\par
\vspace{\questionblockgap}
\noindent\textbf{34.} Computer \uline{graphics}(A) software \uline{has infinite}(B) \uline{applications}(C) in a \uline{widely}(D) array of fields.
\par
\vspace{\questionblockgap}
\noindent\textbf{35.} The planet Mercury \uline{rotates}(A) \uline{slow}(B) than any \uline{other}(C) planet \uline{except}(D) Venus.
\par
\vspace{\questionblockgap}
\noindent\textbf{36.} Van Cliburn \uline{he studied}(A) \uline{piano from}(B) 1951 to 1954 and \uline{won}(C) multiple \uline{awards between}(D) 1958 and 1960.
\par
\vspace{\questionblockgap}
\noindent\textbf{37.} Not \uline{only comics}(A) show a \uline{part of}(B) a story but they \uline{also convey}(C) \uline{the complete}(D) story.
\par
\vspace{\questionblockgap}
\noindent\textbf{38.} How \uline{much information}(A) can \uline{be retained}(B) in short-term memory depends \uline{almost exclusively}(C) on how \uline{it arranged.}(D)
\par
\vspace{\questionblockgap}
\noindent\textbf{39.} When readers \uline{contribute}(A) funds \uline{to private}(B) libraries, these readers \uline{are used}(C) the library \uline{without charge.}(D)
\par
\vspace{\questionblockgap}
\noindent\textbf{40.} \uline{Diagrams}(A) display \uline{informations}(B) in a conspicuous way and \uline{vividly}(C) show differences and \uline{similarities.}(D)
\par
\vspace{\questionblockgap}
\section*{Reading Comprehension}
\begin{center}
\textbf{Time 55 minutes}
\end{center}
DIRECTIONS: In this section you will read several passages. Each passage is followed by a series of questions. For questions 1-50, you need to select the best answer, (A), (B), (C), or (D), to each question. Then, on your answer sheet, find the number of the question and fill in the space that contains the letter of the answer you have selected. Fill in the space completely. Answer all questions following a passage on the basis of what is stated or implied in the passage.\par
\begin{exampleblock}
\small
\sloppy
Read the following passage:\par
A tomahawk is a small ax used as a tool and a weapon by the North American Indian\par
tribes. An average tomahawk was not very long and did not weigh a great deal. Originally,\par
the head of the tomahawk was made of a shaped stone or an animal bone and was mounted on\par
Line a wooden handle. After the arrival of the European settlers, the Indians began to use toma-\par
hawks with iron heads. Indian males and females of all ages used tomahawks to chop and cut wood, pound stakes into the ground to put up wigwams, and do many other chores. Indian warriors relied on tomahawks as weapons and even threw them at their enemies. Some types of tomahawks were used in religious ceremonies. Contemporary American idioms reflect\par
this aspect of American heritage.\par
\end{exampleblock}
\begin{exampleblock}
\textbf{EXAMPLE I}\par
\noindent Early tomahawk heads were made of
\begin{enumerate}[label=(\Alph*),leftmargin=*,itemsep=0.2em]
\item stone or bone
\item wood or sticks
\item European iron
\item religious weapons
\end{enumerate}
\vspace{0.25\baselineskip}
According to the passage, early tomahawk heads were made of stone or bone. Therefore, the correct answer is (A).\par
\end{exampleblock}
\begin{exampleblock}
\textbf{EXAMPLE II}\par
\noindent How has the Indian use of tomahawks affected American daily life today?
\begin{enumerate}[label=(\Alph*),leftmargin=*,itemsep=0.2em]
\item Tomahawks are still used as weapons.
\item Tomahawks are used as tools for.certain jobs.
\item Contemporary language refers to tomahawks.
\item Indian tribes cherish tomahawks as heirlooms.
\end{enumerate}
\vspace{0.25\baselineskip}
The passage states, "Contemporary American idioms reflect this aspect of American heritage." The correct answer is (C). After you read the directions, begin work on the questions.\par
\end{exampleblock}
\subsection*{Questions 1-12}
\begingroup
\small
\sloppy
\setlength{\parskip}{0pt}
\setlength{\parindent}{0pt}
Charles Lindbergh was born in Detroit, Michigan, in 1902 but was raised on a farm in\par
Minnesota, where his father was elected to the U.S. Congress in 1907. From then on, he\par
spent his boyhood alternately in Washington, D.C., Detroit, and Little Falls, Minnesota.\par
Line Because Lindbergh exhibited exceptional mechanical taient, in 1921 he was admitted to\par
the University of Wisconsin to study engineering. However, the young man was seeking\par
more challenging endeavors, and two years later he became a stunt pilot who performed\par
feats at county fairs and public assemblies. This unusual and dangerous undertaking paid\par
off handsomely in the sense that it allowed him to gain a diverse and well-rounded experi-\par
ence in aeronautics. He particularly delighted in what he called "wing-walking" and para-\par
chute jumping.\par
After a year of training as a military cadet, Lindbergh completed his program at the\par
Brooks and Kelly airfields at the top of his class and earned the rank of captain. Robertson\par
Aircraft Corporation of St. Louis, Missouri, offered him employment as a mail pilot to run\par
the routes between St. Louis and Chicago, and Lindbergh retained his position with the\par
company until 1927. During this period, he set out to win the Raymond B. Orteig prize of\par
\$25,000 to be awarded to the first pilot to fly nonstop from New York to Paris. This ambi-\par
tion would irreversibly change his life and accord him a prominent place in the history of\par
Aviation.\par
Embarking on the greatest adventure of his time, Lindbergh left Roosevelt Field at 7:52\par
A.M. on May 20, 1927, and landed at Le Bourget Field at 5:24 P.M. the next day. Fearing that\par
he would be unknown when he arrived, Lindbergh carried letters of introduction to digni-\par
taries in Paris, but when his plane came to a stop, he was overwhelmed by tremendous wel-\par
coming crowds. He was decorated in France, Great Britain, and Belgium, and President\par
Coolidge sent a specially designated cruiser, the Memphis, to bring him back. His accom-\par
plishments in aeronautics brought him more medals and awards than had ever been re-\par
ceived by any other person in private life.\par
\endgroup
\medskip
\noindent\textbf{1.} Which of the following is the best title for the passage?
\begin{enumerate}[label=(\Alph*),leftmargin=*,itemsep=0.2em]
\item A Benchmark Adventure in Aeronautics
\item The Early Life of Charles Lindbergh
\item Groundbreaking Events in Aviation
\item Charles Lindbergh's Explorations
\end{enumerate}
\par
\vspace{0.25\baselineskip}
\noindent\textbf{2.} According to the passage, Lindbergh did not complete his degree because he
\begin{enumerate}[label=(\Alph*),leftmargin=*,itemsep=0.2em]
\item opted for the life of an exhibition pilot
\item pursued training in the military
\item was seeking a sedentary life-style
\item set out to win recognition
\end{enumerate}
\par
\vspace{0.25\baselineskip}
\noindent\textbf{3.} In line 7, the word "assemblies" is closest in meaning to
\begin{enumerate}[label=(\Alph*),leftmargin=*,itemsep=0.2em]
\item hearings
\item houses
\item gatherings
\item shows
\end{enumerate}
\par
\vspace{0.25\baselineskip}
\noindent\textbf{4.} In line 7 the word "undertaking" refers to
\begin{enumerate}[label=(\Alph*),leftmargin=*,itemsep=0.2em]
\item studying at the university
\item exhibiting mechanical talent
\item seeking challenging endeavors
\item performing feats
\end{enumerate}
\par
\vspace{0.25\baselineskip}
\noindent\textbf{5.} In line 8, the word "handsomely" is closest in meaning to
\begin{enumerate}[label=(\Alph*),leftmargin=*,itemsep=0.2em]
\item honorably
\item handily
\item well
\item in time
\end{enumerate}
\par
\vspace{0.25\baselineskip}
\noindent\textbf{6.} It can be inferred from the passage that. as a military cadet, Lindbergh
\begin{enumerate}[label=(\Alph*),leftmargin=*,itemsep=0.2em]
\item was in top form
\item earned a good salary
\item was the best among students
\item trained with the best students
\end{enumerate}
\par
\vspace{0.25\baselineskip}
\noindent\textbf{7.} The author of the passage implies that Lindbergh's job with Robertson Aircraft Corporation
\begin{enumerate}[label=(\Alph*),leftmargin=*,itemsep=0.2em]
\item required regular intercity flights
\item was not intended as longterm employment
\item required him to perform dangerous flights
\item necessitated his running long distances
\end{enumerate}
\par
\vspace{0.25\baselineskip}
\noindent\textbf{8.} In line 17, the word "irreversibly" is closest in meaning to
\begin{enumerate}[label=(\Alph*),leftmargin=*,itemsep=0.2em]
\item forever
\item formerly
\item irresistibly
\item only
\end{enumerate}
\par
\vspace{0.25\baselineskip}
\noindent\textbf{9.} According to the passage, how old was Lindbergh when he carried out his challenging flight?
\begin{enumerate}[label=(\Alph*),leftmargin=*,itemsep=0.2em]
\item Twenty-one
\item Twenty-three
\item Twenty-four
\item Twenty-five
\end{enumerate}
\par
\vspace{0.25\baselineskip}
\noindent\textbf{10.} The author of the passage implies that Lindbergh did not anticipate becoming a
\begin{enumerate}[label=(\Alph*),leftmargin=*,itemsep=0.2em]
\item pilot
\item celebrity
\item mail carrier
\item army captain
\end{enumerate}
\par
\vspace{0.25\baselineskip}
\noindent\textbf{11.} It can be inferred from the passage that in the early 1920s it was NOT common for young people to
\begin{enumerate}[label=(\Alph*),leftmargin=*,itemsep=0.2em]
\item study engineering
\item train as officers
\item go on exhibition tours
\item be elected to an office
\end{enumerate}
\par
\vspace{0.25\baselineskip}
\noindent\textbf{12.} A paragraph following the passage would most probably discuss
\begin{enumerate}[label=(\Alph*),leftmargin=*,itemsep=0.2em]
\item the development of commercial and military aviation
\item the reaction of the government to Lindbergh's flight
\item the effect of instant celebrity on Lindbergh
\item Lindbergh's aircraft and engine modifications
\end{enumerate}
\par
\vspace{0.25\baselineskip}
\subsection*{Questions 13-21}
\begingroup
\small
\sloppy
\setlength{\parskip}{0pt}
\setlength{\parindent}{0pt}
Lithography is a planographic process that performs a significant function in illustration\par
and offset printing. It is based o n the principle that water does not combine with grease-\par
based substances, preventing them from smearing an outline on an unpolished surface. The\par
Line contour does not need to be engraved into the plate, as in the case of gravure printing, or\par
raised above the surface, as in the letterpress process. These laborious operations ensure that\par
only the design to be printed catches and retains the ink transferred to the paper.\par
In lithography, the artist draws on a leveled, grainy plate made of limestone, zinc, alumi-\par
num, or specially treated paper with a grease pencil, a crayon, or tusche, a greasy liquid. Af\par
ter Sketching the contour on the plate, the artist coats both the drawn and the undrawn por-\par
tions of the plate with an inking roller dipped in a solution of nitric acid and gum arabic.\par
The gum arabic envelops the greased surfaces and prevents ink from penetrating into the\par
greaseless areas. The artist dampens the surface with water, which is repelled by the greased\par
areas. Then the surface is covered with thick, oily ink and pressed onto paper. The sheet\par
picks up the ink from the design while the damp stone around the patter keeps the ink\par
from spreading.\par
In offset lithography, shiny sheets of zinc and aluminum are used instead of the heavy,\par
hard-to-handle stone plates. The metal plates are scoured by emery dust and marble chips to\par
give them a grained finish. The subjects to be printed are laid down photographically, and\par
rotary presses automatically moisten, ink, and print hundreds of impressions per hour.\par
\endgroup
\medskip
\noindent\textbf{13.} What does the passage mainly discuss?
\begin{enumerate}[label=(\Alph*),leftmargin=*,itemsep=0.2em]
\item Commercial printing of mass-produced lithographs
\item Steps in a technique for making impressions
\item The equipment necessary for offset lithographs
\item The evolution of lithograph printing to rotary presses
\end{enumerate}
\par
\vspace{0.25\baselineskip}
\noindent\textbf{14.} According to the passage, lithographic printing makes use of the fact that
\begin{enumerate}[label=(\Alph*),leftmargin=*,itemsep=0.2em]
\item artists can draw on flat, greaseless surfaces
\item oily substances do not mix with water
\item gravure etching is work- and time-consuming
\item limestone, zinc, and aluminum can be used as planes
\end{enumerate}
\par
\vspace{0.25\baselineskip}
\noindent\textbf{15.} In line 4, the word "contour" is closest in meaning to
\begin{enumerate}[label=(\Alph*),leftmargin=*,itemsep=0.2em]
\item contrast
\item content
\item outline
\item edge
\end{enumerate}
\par
\vspace{0.25\baselineskip}
\noindent\textbf{16.} In line 7, the word "leveled" is closest in meaning to
\begin{enumerate}[label=(\Alph*),leftmargin=*,itemsep=0.2em]
\item elevated
\item low
\item flawed
\item flattened
\end{enumerate}
\par
\vspace{0.25\baselineskip}
\noindent\textbf{17.} In line 11, the word "envelops" i s closest in meaning to
\begin{enumerate}[label=(\Alph*),leftmargin=*,itemsep=0.2em]
\item sends
\item soils
\item coats
\item coils
\end{enumerate}
\par
\vspace{0.25\baselineskip}
\noindent\textbf{18.} It can be inferred from the passage that in making lithographic prints, the paper
\begin{enumerate}[label=(\Alph*),leftmargin=*,itemsep=0.2em]
\item absorbs the ink from the printing plate
\item spreads the ink on the greased areas
\item shrinks away from the printing stone
\item keeps the oil from sliding off
\end{enumerate}
\par
\vspace{0.25\baselineskip}
\noindent\textbf{19.} Where in the passage does the author point out the advantages of lithography over other types of printing?
\begin{enumerate}[label=(\Alph*),leftmargin=*,itemsep=0.2em]
\item Lines 1-2
\item Lines 3-5
\item Lines 7-8
\item Lines 16-18
\end{enumerate}
\par
\vspace{0.25\baselineskip}
\noindent\textbf{20.} A paragraph following the passage would most probably discuss
\begin{enumerate}[label=(\Alph*),leftmargin=*,itemsep=0.2em]
\item photosynthesis in commercial lithographs
\item the offset printing of billboard advertisements
\item technological advancements in offset printing
\item types of unique lithographs. printed in rare books
\end{enumerate}
\par
\vspace{0.25\baselineskip}
\noindent\textbf{21.} In line 19, the word "impressions" is closest in meaning to
\begin{enumerate}[label=(\Alph*),leftmargin=*,itemsep=0.2em]
\item originals
\item reproductions
\item photographs
\item plates
\end{enumerate}
\par
\vspace{0.25\baselineskip}
\subsection*{Questions 22-32}
\begingroup
\small
\sloppy
\setlength{\parskip}{0pt}
\setlength{\parindent}{0pt}
Although a few protozoans are multicellular, the simplest are unicellular organisms,\par
Such as amoebas, bacteria, sarcodina, ciliates, flagellates, and sporozoans, which can be\par
amorphous in shape and smaller than.001 inch. Cytoplasm fills the cell membrane that en-\par
Line closes it and functions as a barrier between cells. The membrane serves as the outer tissue,\par
and any compound that may destroy the cell has to penetrate it to reach the cytoplasm.\par
Some types of organisms are termed colonial because they represent loosely assembled\par
groups of structurally similar and unifunctional cells. Colonial organisms maintain a\par
symbiotic relationship within their particular environments.\par
Unlike colonial organisms, almost all species of animals and plants are multicellular and\par
include various types of specialized or somatic cells, each with its own nucleus, genetic code,\par
and RNA. The overall size of a multicellular body is contingent on the total number of cells\par
that comprise it, not the size of individual cells. The simplest multicellular animals are hy-\par
dras, sponges, and jellyfish, which have well-defined tissues, a cellular nucleus, and an ele-\par
ment of cell functions. Sponges have a few specialized cells but largely resemble colonial or\par
ganisms that can readily form a new individual group. If the cells of a sponge are separated,\par
they rejoin and continue as a newly formed colonial organism.\par
\endgroup
\medskip
\noindent\textbf{22.} How many cells do the simplest organisms contain?
\begin{enumerate}[label=(\Alph*),leftmargin=*,itemsep=0.2em]
\item One
\item One hundred
\item Many
\item An unknown number
\end{enumerate}
\par
\vspace{0.25\baselineskip}
\noindent\textbf{23.} In line 4, the phrase "outer tissue" is closest in meaning to
\begin{enumerate}[label=(\Alph*),leftmargin=*,itemsep=0.2em]
\item outside force
\item outlying area
\item shell
\item cell
\end{enumerate}
\par
\vspace{0.25\baselineskip}
\noindent\textbf{24.} In line 5, the word "it" refers to
\begin{enumerate}[label=(\Alph*),leftmargin=*,itemsep=0.2em]
\item the cell
\item any compound
\item the membrane
\item the cytoplasm
\end{enumerate}
\par
\vspace{0.25\baselineskip}
\noindent\textbf{25.} It can be inferred from the passage that a cell serves as
\begin{enumerate}[label=(\Alph*),leftmargin=*,itemsep=0.2em]
\item a partition of organism functions
\item the smallest colonial group
\item the smallest genetic unit
\item a flagellate reproductive organ
\end{enumerate}
\par
\vspace{0.25\baselineskip}
\noindent\textbf{26.} In line 6, the word "loosely" is closest in meaning to
\begin{enumerate}[label=(\Alph*),leftmargin=*,itemsep=0.2em]
\item lively
\item naturally
\item freely
\item feebly
\end{enumerate}
\par
\vspace{0.25\baselineskip}
\noindent\textbf{27.} In line 8, the word "symbiotic" is closest in meaning to
\begin{enumerate}[label=(\Alph*),leftmargin=*,itemsep=0.2em]
\item mutually dependent
\item mutually exclusive
\item mutually hostile
\item mutually resistant
\end{enumerate}
\par
\vspace{0.25\baselineskip}
\noindent\textbf{28.} The author of the passage implies that large animals and plants have
\begin{enumerate}[label=(\Alph*),leftmargin=*,itemsep=0.2em]
\item larger cell sizes than amoebas and protozoans
\item larger quantities of protoplasm than smaller lifeforms
\item stronger cellular membranes than flagellates
\item a greater number of cells than smaller lifeforms
\end{enumerate}
\par
\vspace{0.25\baselineskip}
\noindent\textbf{29.} In line 10, the word "each" refers to
\begin{enumerate}[label=(\Alph*),leftmargin=*,itemsep=0.2em]
\item animals
\item species
\item cells
\item plants
\end{enumerate}
\par
\vspace{0.25\baselineskip}
\noindent\textbf{30.} According to the passage, sponges and jellytish are
\begin{enumerate}[label=(\Alph*),leftmargin=*,itemsep=0.2em]
\item the simplest one-cell organisms
\item the simplest multicellular animals
\item tissues and cellular nuclei
\item cellular and colonial organisms
\end{enumerate}
\par
\vspace{0.25\baselineskip}
\noindent\textbf{31.} The author mentions all of the following EXCEPT:
\begin{enumerate}[label=(\Alph*),leftmargin=*,itemsep=0.2em]
\item procreative mechanisms
\item colonial organisms
\item cell contents
\item specialized cells
\end{enumerate}
\par
\vspace{0.25\baselineskip}
\noindent\textbf{32.} This passage would most likely be found in a textbook on which of the following subjects?
\begin{enumerate}[label=(\Alph*),leftmargin=*,itemsep=0.2em]
\item Genetics
\item Anatomy
\item Biology
\item Biochemistry
\end{enumerate}
\par
\vspace{0.25\baselineskip}
\subsection*{Questions 33-42}
\begingroup
\small
\sloppy
\setlength{\parskip}{0pt}
\setlength{\parindent}{0pt}
When parchment, which was extraordinarily costly, was replaced by papyrus, it became\par
feasible to establish libraries. At the onset, they began as archives for record keeping and\par
document storage. According to second-hand reports, the most renowned papyrus library\par
Line was the Alexandrian, founded by Alexander the Great around 330 B.C. in Alexandria, Egypt.\par
His successors as rulers of Egypt, Ptolemy I and Ptolemy II, expanded the library into the\par
greatest collection of scrolls in the ancient world. To acquire this collection, the rulers bor-\par
rowed scrolls and manuscripts from libraries in Athens, Rome, and other localities and or-\par
dered them duplicated. At times, the library employed more than 100 scribes and illustra-\par
tors. Some historians claim that the Alexandrian library purchased entire lesser libraries to\par
contribute to and enhance the quality of its possessions.\par
The library owned a copy of every contemporary scroll known to the library's administra-\par
tors and contained more than 400,000 items, all of which were classified and organized. The\par
contents of the papyrus rolls were edited, and a bibliography of Greek literature was com-\par
piled and cross-referenced, reflecting the emergence and dissemination of a highly devel-\par
oped Greek culture. Over time, a succession of leading scholars directed this library, which\par
was acclaimed for the scholarly undertakings it supported as well as for the size of its collec-\par
tion. At one time, 72 scholars were engaged to translate religious testaments, historical an-\par
nals, and mercantile accounts. Although the library flourished, it was accessible to only a\par
minority of the population because in ancient times the vast majority of urban dwellers were.\par
illiterate. Because papyrus was extremely perishable, not a trace of the Alexandrian library\par
remains today, and archaeologists have several hypotheses as to what became of it.\par
\endgroup
\medskip
\noindent\textbf{33.} What does the passage mainly discuss?
\begin{enumerate}[label=(\Alph*),leftmargin=*,itemsep=0.2em]
\item The use of papyrus in ancient scroll collections
\item The origin an d history of a library
\item The cultural initiatives of Alexander the Great
\item The expansion of libraries in ancient times
\end{enumerate}
\par
\vspace{0.25\baselineskip}
\noindent\textbf{34.} In line 2, the word "feasible" is closest in meaning to
\begin{enumerate}[label=(\Alph*),leftmargin=*,itemsep=0.2em]
\item practicable
\item easy
\item prestigious
\item ebullient
\end{enumerate}
\par
\vspace{0.25\baselineskip}
\noindent\textbf{35.} It can be inferred from the passage that reports of the Alexandrian library
\begin{enumerate}[label=(\Alph*),leftmargin=*,itemsep=0.2em]
\item were highly exaggerated
\item could not be verified
\item were secondary in importance
\item could not be made known
\end{enumerate}
\par
\vspace{0.25\baselineskip}
\noindent\textbf{36.} The author of the passage implies that the rulers of Egypt
\begin{enumerate}[label=(\Alph*),leftmargin=*,itemsep=0.2em]
\item oversaw the expansion of the library directly
\item devoted funds and other resources to the library collections
\item sought to make the library self-contained
\item marshaled worldwide support for the library collections
\end{enumerate}
\par
\vspace{0.25\baselineskip}
\noindent\textbf{37.} According to the passage, the main goal of the library in Alexandria was
\begin{enumerate}[label=(\Alph*),leftmargin=*,itemsep=0.2em]
\item collecting scrolls loaned by other libraries
\item gradually replacing papyrus with parchment
\item translating scrolls in ancient Egypt and Greece
\item accumulating translations and originals of texts
\end{enumerate}
\par
\vspace{0.25\baselineskip}
\noindent\textbf{38.} In the second paragraph, the author. implies that
\begin{enumerate}[label=(\Alph*),leftmargin=*,itemsep=0.2em]
\item parchment was more durable than books
\item libraries were necessary to conduct research
\item the library collection cannot be examined
\item the library was historically relevant
\end{enumerate}
\par
\vspace{0.25\baselineskip}
\noindent\textbf{39.} With which of the following statements about Greek literature is the author of the passage most likely to agree?
\begin{enumerate}[label=(\Alph*),leftmargin=*,itemsep=0.2em]
\item It was nurtured in libraries in Athens and Rome.
\item It was integral to Greek culture.
\item It was compiled and cross-referenced in the library.
\item It was beginning to emerge when the library was expanded.
\end{enumerate}
\par
\vspace{0.25\baselineskip}
\noindent\textbf{40.} In line 15, the word "succession" is closest in meaning to
\begin{enumerate}[label=(\Alph*),leftmargin=*,itemsep=0.2em]
\item series
\item success
\item sundry
\item substitution
\end{enumerate}
\par
\vspace{0.25\baselineskip}
\noindent\textbf{41.} It can be inferred from the passage that in ancient times
\begin{enumerate}[label=(\Alph*),leftmargin=*,itemsep=0.2em]
\item books and scrolls were updated regularly
\item libraries benefited upper social classes
\item maintaining collections was fruitless
\item the population should have been educated
\end{enumerate}
\par
\vspace{0.25\baselineskip}
\noindent\textbf{42.} In the last sentence, the phrase "not a trace" most probably means
\begin{enumerate}[label=(\Alph*),leftmargin=*,itemsep=0.2em]
\item absolutely no one
\item absolutely nothing
\item not a penny
\item not a soul
\end{enumerate}
\par
\vspace{0.25\baselineskip}
\subsection*{Questions 43-50}
\begingroup
\small
\sloppy
\setlength{\parskip}{0pt}
\setlength{\parindent}{0pt}
According to data obtained from radioactive dating, the oldest rocks found on earth are\par
approximately 500 million to 4 billion years old. Similar ages have been determined for me-\par
teorites and the rocks gathered from the moon's surface. Different methods of arriving at the\par
Line earth's age generate very similar results. Modern theories about the formation, develop-\par
ement, and eventual burning out of stars suggest that the sun is about 5 billion years old. Ex-\par
perts contend that the earth and the sun were formed at almost the same time from a cloud of\par
dust and gas resulting from a cosmic explosion. The present rate of expansion of the galaxies\par
can be extrapolated to suggest that, if the universe began with a "big bang" about 15 billion\par
years ago, an age of 5 billion years for both the earth and the sun can be considered plausible.\par
Long before radioactive dating was implemented, mythology and oral narratives alluded\par
to a conjecture that the earth was nearly 6,000 years old. The methods of computation based\par
on the analysis of genealogical trees in scant archaeological findings provide evidence that\par
can be difficult to date accurately. Today, radioactive dating of particles and whole objects\par
has rejected this figure of the earth's age as unreliable.\par
\endgroup
\medskip
\noindent\textbf{43.} What does the passage mainly discuss?
\begin{enumerate}[label=(\Alph*),leftmargin=*,itemsep=0.2em]
\item Dating techniques in research
\item Modern theories and radioactive dating
\item Research and narratives about the earth's formation
\item Establishing the earth's age
\end{enumerate}
\par
\vspace{0.25\baselineskip}
\noindent\textbf{44.} It can be inferred from the passage that radioactive dating is important for estimating the age of
\begin{enumerate}[label=(\Alph*),leftmargin=*,itemsep=0.2em]
\item all known meteors
\item all existing planets
\item the earth
\item the trees
\end{enumerate}
\par
\vspace{0.25\baselineskip}
\noindent\textbf{45.} In line 5, the word "eyentual" is closest in meaning to
\begin{enumerate}[label=(\Alph*),leftmargin=*,itemsep=0.2em]
\item ultimate
\item eventful
\item utter
\item enduring
\end{enumerate}
\par
\vspace{0.25\baselineskip}
\noindent\textbf{46.} According to the passage, the moon is
\begin{enumerate}[label=(\Alph*),leftmargin=*,itemsep=0.2em]
\item older than the earth and the sun
\item newer than the earth and the sun
\item approximately the same age as the earth and the sun
\item approximately the same density as the earth and sun
\end{enumerate}
\par
\vspace{0.25\baselineskip}
\noindent\textbf{47.} The author of the passage implies that
\begin{enumerate}[label=(\Alph*),leftmargin=*,itemsep=0.2em]
\item the earth and the sun are of similar origin
\item the earth and the sun can be explosive
\item meteorites and the moon have been analyzed
\item the galaxies are expanding at a substantial rate
\end{enumerate}
\par
\vspace{0.25\baselineskip}
\noindent\textbf{48.} With which of the following statements would the author be most likely to agree?
\begin{enumerate}[label=(\Alph*),leftmargin=*,itemsep=0.2em]
\item The moon and the sun are 15 billion years old.
\item The moon can be viewed as a meteorite.
\item The formation of galaxies is an on going process.
\item The earth can be dated a s far back as 6,000 years.
\end{enumerate}
\par
\vspace{0.25\baselineskip}
\noindent\textbf{49.} In line 9, the word "plausible" is closest in meaning to
\begin{enumerate}[label=(\Alph*),leftmargin=*,itemsep=0.2em]
\item reasonable
\item rational
\item relative
\item relational
\end{enumerate}
\par
\vspace{0.25\baselineskip}
\noindent\textbf{50.} What conclusion does the author of the passage make?
\begin{enumerate}[label=(\Alph*),leftmargin=*,itemsep=0.2em]
\item Radioactive dating is refuted by researchers.
\item Radioactive dating is more accurate than other methods.
\item The earth is a part of a galaxy that includes many moons.
\item The sun's radioactivity is scant and can be negligible.
\end{enumerate}
\par
\vspace{0.25\baselineskip}
\end{document}
